\documentclass[10pt]{scrartcl}
\usepackage[mdthm,colorsec,noauthor,noasy]{evan}
\usepackage{handout}
\DeclareMathOperator{\im}{range}
\title{MATH 110 --- Week 3}
\subtitle{Notes Chapters 5--6; Homework 3}
\author{}\date{}\ihead{}
\begin{document}\maketitle
\section{Linear maps and their operations}\begin{qbox}Which functions preserve linear combinations? The map $T(x,y)=(2x+y,x)$ does, since each output is a linear combination of the inputs.\end{qbox}
\begin{prob}\begin{definition*}[Notes Definition 11; LADR 3.1]
Let $U,V,W$ be vector spaces over the same field $F$. A map $T:V\to W$ is linear when it is additive, $T(v+w)=T(v)+T(w)$, and homogeneous, $T(av)=aT(v)$, for all $a\in F$ and $v,w\in V$.
\end{definition*}\end{prob}
\begin{proposition*}[Notes Definition 11; LADR 3.10]
A linear map satisfies $T0=0$ and preserves every finite linear combination. Conversely, $T(av+w)=aTv+Tw$ for all $a,v,w$ suffices for linearity.
\end{proposition*}
\begin{motive}\ii[Want.] Use one combination test and preserve zero.\ii[Know.] Vector addition permits cancellation.\ii[Key step.] Start with zero and then build a finite sum one term at a time.\end{motive}
\begin{prob}\begin{proof}\begin{steps}
\ii Additivity gives $T0=T(0+0)=T0+T0$; cancel $T0$ to obtain $T0=0$.
\ii Homogeneity handles one term. Inductively apply additivity to a sum of the first $k$ terms and the last term, obtaining $T(\sum_j a_jv_j)=\sum_j a_jTv_j$.
\ii Conversely the combination test with $a=1$ gives additivity, hence zero preservation as above; then $w=0$ gives homogeneity.
\end{steps}\end{proof}\end{prob}
\begin{example*}[Notes Example 21]For $n\ge1$, differentiation $D:\mathcal P_n(\mathbb R)\to\mathcal P_{n-1}(\mathbb R)$ is linear by the power rule and distributivity. We cite the elementary calculus facts that differentiation and integration are linear and that $(xf)'=f+xf'$.
\begin{steps}\ii On the complex vector space of continuous functions satisfying $f(t+2\pi)=f(t)$, define $c_k(f)=\frac1{2\pi}\int_0^{2\pi}f(t)e^{-ikt}\,dt$ for each integer $k$. Here $e^{-ikt}=\cos(kt)-i\sin(kt)$. These are the Fourier coefficients; the integral of a complex function means the separate real and imaginary integrals.
\ii Integral linearity gives $c_k(af+bg)=ac_k(f)+bc_k(g)$ for every $k$, so $f\mapsto(c_k(f))_{k\in\mathbb Z}$ is linear.\end{steps}\end{example*}
\begin{proposition*}[Notes Proposition 18; LADR 3.6]$L(V,W)$, the set of linear maps $V\to W$, is a vector space under $(S+T)(v)=S(v)+T(v)$ and $(aT)(v)=aT(v)$.\end{proposition*}
\begin{motive}\ii[Want.] Check that combining maps preserves linearity.\ii[Know.] The space $W^V$ of all functions has pointwise vector operations (Week 1).\ii[Key step.] Apply the subspace test inside $W^V$.\end{motive}
\begin{prob}\begin{proof}\begin{steps}\ii The zero map is linear. For linear $S,T$, evaluate $(aS+T)(bv+w)=a(bS(v)+S(w))+bT(v)+T(w)$.\ii Regrouping gives $b(aS+T)(v)+(aS+T)(w)$. Thus $aS+T$ is linear, proving the subspace condition.\end{steps}\end{proof}\end{prob}
\begin{prob}\begin{definition*}[Notes Definition 12; LADR 5.1]An operator is a linear map $V\to V$; write $L(V)=L(V,V)$. The identity is $I_V(v)=v$, and the zero operator sends every vector to zero. Composition is $ST=S\circ T$, so $T$ acts first. For an operator, $T^0=I_V$ and $T^{k+1}=TT^k$.\end{definition*}\end{prob}
\begin{proposition*}[Notes Proposition 19; LADR 3.7--3.8]If $T:U\to V$ and $S:V\to W$ are linear, then $ST$ is linear. Composition is associative, has identities $TI_U=I_VT=T$, and distributes on both sides over sums of compatible maps.\end{proposition*}
\begin{motive}\ii[Want.] Compute with compositions as algebraic products.\ii[Know.] Equality of maps can be checked at each input.\ii[Key step.] Expand at a general vector, respecting the order of application.\end{motive}
\begin{prob}\begin{proof}\begin{steps}\ii $ST(au+v)=S(aTu+Tv)=aSTu+STv$, proving linearity.\ii For compatible $R,S,T$, $((RS)T)u=R(S(Tu))=(R(ST))u$. Also $(TI_U)u=Tu=(I_VT)u$.\ii $((S_1+S_2)T)u=S_1Tu+S_2Tu$ and $(S(T_1+T_2))u=ST_1u+ST_2u$ by linearity of $S$.\end{steps}\end{proof}\end{prob}
\begin{example*}[Notes Examples 22--23]\begin{steps}\ii On real sequences indexed by $0,1,\ldots$, the shift $S(a_0,a_1,\ldots)=(a_1,a_2,\ldots)$ is linear coordinatewise. If $a_{n+2}=a_{n+1}+a_n$, its shifted sequence satisfies the same relation; hence on this recurrence subspace $S^2-S-I=0$ by comparison of every coordinate.
\ii On $\mathbb R^2$, $P(x,y)=(0,y)$, $R_\theta(x,y)=(x\cos\theta-y\sin\theta,x\sin\theta+y\cos\theta)$, and $N(x,y)=(y,0)$ are linear by distributing each coordinate.
\ii $R_{-\pi/2}P(x,y)=(y,0)=N(x,y)$ and $N^2(x,y)=N(y,0)=0$. On the standard basis, $e_2\mapsto e_1\mapsto0$.\end{steps}
\begin{center}\begin{tikzpicture}[>=stealth]\node (a) at (0,0) {$e_2$};\node (b) at (2,0) {$e_1$};\node (c) at (4,0) {$0$};\draw[->] (a)--node[above]{$N$}(b);\draw[->] (b)--node[above]{$N$}(c);\end{tikzpicture}\\The second application of $N$ sends both standard basis vectors to zero.\end{center}\end{example*}
\begin{example*}[Notes Example 24]A smooth function is one with derivatives of every order. On smooth real functions define $Xf(x)=xf(x)$ and $Df=f'$. These are linear and again produce smooth functions by elementary calculus.
\begin{steps}\ii The product rule gives $(DX)f=(xf)'=f+xf'$, whereas $(XD)f=xf'$.
\ii Therefore $DX-XD=I$. This is nonzero because $I(1)=1$; thus two operators need not commute.\end{steps}\end{example*}
\begin{shortcuts}\ii Check $T(av+w)=aTv+Tw$ once; it gives both linearity conditions.\ii To disprove linearity, first test $T(0)=0$; every linear map satisfies it.\end{shortcuts}\begin{trapbox}\ii $T(x)=x+1$ is not linear since $T(0)=1$. Passing the zero test alone is insufficient: $T(x)=x^2$ sends zero to zero but $T(2)=4\ne2T(1)$.\ii Composition order matters: for the maps above, $PN=0$ while $NP=N\ne0$.\end{trapbox}
\section{Constructing a map from specified values}\begin{qbox}Can images of basis vectors be chosen freely? On $\mathbb R^2$, specifying $Te_1=u$ and $Te_2=w$ forces $T(x,y)=xu+yw$.\end{qbox}
\begin{prob}\begin{definition*}[Notes Definition 13]For a subspace $U\subseteq V$, the restriction $T|_U:U\to W$ has the same rule $T|_U(u)=T(u)$ and is linear because $au+v\in U$ and $T$ is linear.\end{definition*}\end{prob}
\begin{proposition*}[Notes Proposition 20]For a direct sum $U_1\oplus\cdots\oplus U_m$ and linear maps $T_j:U_j\to W$, there is a unique linear map $T$ on the sum whose restriction to each $U_j$ is $T_j$.\end{proposition*}
\begin{motive}\ii[Want.] Combine the specified maps into one map.\ii[Know.] Week 1 gives a unique decomposition into direct-sum components.\ii[Key step.] Define the output by adding the outputs of the unique components.\end{motive}
\begin{prob}\begin{proof}\begin{steps}\ii Define $T(u_1+\cdots+u_m)=\sum_jT_j(u_j)$. Unique components make this well-defined.\ii The components of $a\sum u_j+\sum v_j$ are $au_j+v_j$, so their images sum to $a\sum T_ju_j+\sum T_jv_j$. This proves linearity.\ii If only the $j$th component is nonzero, the rule agrees with $T_j$. Any linear extension must add those prescribed component images, proving uniqueness.\end{steps}\end{proof}\end{prob}
\begin{proposition*}[Notes Proposition 21; LADR 3.4]For a basis $v_1,\ldots,v_n$ of $V$ and arbitrary $w_1,\ldots,w_n\in W$, there is a unique $T\in L(V,W)$ with $Tv_j=w_j$.\end{proposition*}
\begin{motive}\ii[Want.] Build a map with prescribed basis images.\ii[Know.] Week 2 gives unique basis coefficients, and coefficients add and scale coordinatewise.\ii[Key step.] Apply the same coefficients to the chosen images.\end{motive}
\begin{prob}\begin{proof}\begin{steps}\ii Define $T(\sum a_jv_j)=\sum a_jw_j$; uniqueness of the $a_j$ makes it well-defined.\ii Expanding two inputs in the basis shows that $T(au+v)=aTu+Tv$. Substitution of $v_j$ gives $Tv_j=w_j$.\ii Any linear map with these basis values must have the displayed rule, proving uniqueness.\end{steps}\end{proof}\end{prob}
\method{Prescribing images on a nonstandard basis}\begin{steps}\ii In $\mathbb R^2$, let $v_1=(1,1),v_2=(1,-1)$; solving $(x,y)=av_1+bv_2$ gives $a=(x+y)/2,b=(x-y)/2$. Thus the list is a basis.\ii Prescribe $Tv_1=(1,0)$ and $Tv_2=(0,2)$. Proposition 21 gives $T(x,y)=((x+y)/2,x-y)$.\ii Substituting the two basis vectors checks the prescribed values.\end{steps}
\begin{shortcuts}\ii Two linear maps agreeing on a basis agree everywhere by Proposition 21.\ii To define a map on a direct sum, define it independently on each summand; unique decomposition makes the rules compatible.\end{shortcuts}\begin{trapbox}\ii Images of a dependent list cannot be freely prescribed: if $v_3=v_1+v_2$, linearity forces $Tv_3=Tv_1+Tv_2$.\ii A non-direct sum requires compatibility: assigning different values to the same nonzero vector in an intersection cannot define a function.\end{trapbox}
\section{Kernel, range and rank--nullity}\begin{qbox}How much freedom remains after applying a linear map? For $P(x,y)=(0,y)$, the first coordinate is lost and the second is retained.\end{qbox}
\begin{prob}\begin{definition*}[Notes Definition 14; LADR 3.11, 3.16]The kernel (null space) is $\ker T=\{v\in V:Tv=0\}$; the range (image) is $\im T=\{Tv:v\in V\}\subseteq W$.\end{definition*}\end{prob}
\begin{proposition*}[Notes Proposition 22; LADR 3.13, 3.18]For $T\in L(V,W)$, $\ker T$ is a subspace of $V$ and $\im T$ is a subspace of $W$.\end{proposition*}
\begin{motive}\ii[Want.] Use dimension and bases for the two sets.\ii[Know.] We proved $T0=0$ and know the subspace test.\ii[Key step.] Express a range vector as an image before combining it.\end{motive}
\begin{prob}\begin{proof}\begin{steps}\ii Zero belongs to both sets. If $u,v\in\ker T$, then $T(au+v)=aTu+Tv=0$.\ii If $w_1=Tv_1$ and $w_2=Tv_2$, then $aw_1+w_2=T(av_1+v_2)$ lies in the range. The subspace test proves both claims.\end{steps}\end{proof}\end{prob}
\begin{prob}\begin{definition*}[Notes Definition 15]The dimensions of $\ker T$ and $\im T$ are called the nullity and rank of $T$, respectively.\end{definition*}\end{prob}
\begin{example*}[Notes Example 25]For $D\in L(\mathcal P_4(\mathbb R))$, $Dp=p'$.
\begin{steps}\ii If $p=a_0+a_1x+\cdots+a_4x^4$, then $Dp=a_1+2a_2x+3a_3x^2+4a_4x^3$. Coefficient comparison gives $\ker D=\operatorname{span}(1)$.
\ii The images of the monomial basis are $0,1,2x,3x^2,4x^3$, which span $\mathcal P_3(\mathbb R)$. Every image has degree at most $3$, so $\im D=\mathcal P_3(\mathbb R)$, nullity $1$, rank $4$.\end{steps}\end{example*}
\begin{proposition*}[Notes Proposition 23; LADR 3.15, 3.19]$T$ is injective iff $\ker T=\{0\}$; it is surjective iff $\im T=W$.\end{proposition*}
\begin{motive}\ii[Want.] Turn equality of outputs into a condition on one vector.\ii[Know.] Linearity gives $Tu-Tv=T(u-v)$.\ii[Key step.] Apply the kernel condition to the difference of two inputs.\end{motive}
\begin{prob}\begin{proof}\begin{steps}\ii If $T$ is injective and $Tv=0=T0$, then $v=0$.\ii If $\ker T=\{0\}$ and $Tu=Tv$, then $u-v\in\ker T$, so $u=v$.\ii Surjectivity means every element of $W$ is $Tv$ for some $v$, exactly $\im T=W$.\end{steps}\end{proof}\end{prob}
\begin{proposition*}[Notes Proposition 24; LADR 3.21]If $V$ is finite-dimensional and $T:V\to W$ is linear, a basis $u_1,\ldots,u_m$ of $\ker T$ extends to a basis $u_1,\ldots,u_m,v_1,\ldots,v_n$ of $V$ such that $Tv_1,\ldots,Tv_n$ is a basis of $\im T$. Hence $\dim V=\dim\ker T+\dim\im T$.\end{proposition*}
\begin{motive}\ii[Want.] Count exactly the directions whose images remain independent.\ii[Know.] Week 2 allows a kernel basis to be extended to a domain basis.\ii[Key step.] Apply $T$ to the added vectors; relations among their images become relations in the domain basis.\end{motive}
\begin{prob}\begin{proof}\begin{steps}\ii The kernel is a finite-dimensional subspace, so choose its basis and extend it as stated by Week 2.\ii If $\sum a_jTv_j=0$, then $\sum a_jv_j\in\ker T$, so it equals $\sum b_iu_i$. Independence of the full domain basis forces all $a_j,b_i=0$.\ii For any $w=Tv$, expand $v=\sum b_iu_i+\sum a_jv_j$; then $w=\sum a_jTv_j$. This proves spanning.\ii The dimensions are respectively $m+n,m,n$, giving the formula. No finite-dimensional hypothesis on $W$ is needed.\end{steps}\end{proof}\end{prob}
\begin{proposition*}[Notes Proposition 25; LADR 3.65]If $\dim V=\dim W<\infty$, a linear map $T:V\to W$ is injective iff surjective iff bijective.\end{proposition*}
\begin{motive}\ii[Want.] Deduce one bijectivity condition from the other.\ii[Know.] Rank--nullity and the fact that an equal-dimensional subspace is the full space (Week 2).\ii[Key step.] Compare the rank with the common dimension.\end{motive}
\begin{prob}\begin{proof}\begin{steps}\ii If $T$ is injective, its kernel is zero, so rank--nullity gives $\dim\im T=\dim V=\dim W$. Thus $\im T=W$.\ii If $T$ is surjective, its rank is $\dim W=\dim V$, so its nullity is zero. A zero-dimensional space is $\{0\}$, proving injectivity.\ii Bijective means both properties, completing the equivalence.\end{steps}\end{proof}\end{prob}
\method{Compute kernel and range from a formula}\begin{steps}\ii For $P(x,y)=(0,y)$, $Pv=0$ means $y=0$, so $\ker P=\mathbb Re_1$; outputs are exactly $\mathbb Re_2$. These are complementary.\ii For $N(x,y)=(y,0)$, the same equation $y=0$ gives $\ker N=\mathbb Re_1$, but now $\im N=\mathbb Re_1$ too.\ii Both satisfy $2=1+1$, yet only the first gives a direct-sum decomposition.\end{steps}
\begin{proposition*}[LADR 3.22, 3.24, 3.26, 3.28]For finite-dimensional $V,W$: if $\dim V>\dim W$, no linear $T:V\to W$ is injective; if $\dim V<\dim W$, none is surjective. A homogeneous system has zero right-hand sides. Thus a homogeneous linear system with more unknowns than equations has a nonzero solution, while a system with more equations than unknowns has some right-hand sides with no solution.\end{proposition*}
\begin{motive}\ii[Want.] Decide impossibility from dimensions before computing.\ii[Know.] Rank--nullity and $\dim\im T\le\dim W$.\ii[Key step.] Bound the rank from the domain and codomain separately.\end{motive}
\begin{prob}\begin{proof}\begin{steps}\ii Rank--nullity gives $\dim\ker T=\dim V-\dim\im T\ge\dim V-\dim W$. If the last quantity is positive, the kernel is nonzero.\ii Also $\dim\im T=\dim V-\dim\ker T\le\dim V$. If $\dim V<\dim W$, the range cannot be all of $W$.\ii A system of $r$ equations in $s$ unknowns defines a linear map $F^s\to F^r$ by its left sides. Apply the first bound if $s>r$ and the second if $s<r$.\end{steps}\end{proof}\end{prob}
\begin{shortcuts}\ii To find a range, apply $T$ to any spanning list: linearity says the images span the range.\ii For equal finite dimensions, a zero kernel proves surjectivity without solving every output equation.\end{shortcuts}\begin{trapbox}\ii Rank--nullity is a dimension identity, not $V=\ker T\oplus\im T$: the map $N$ is a counterexample.\ii Equal dimension is essential: $x\mapsto(x,0)$ from $\mathbb R$ to $\mathbb R^2$ is injective but not onto. Finite dimension is essential too: the sequence shift is onto but has nonzero kernel.\end{trapbox}
\section{Coordinates and representation matrices}\begin{qbox}How can a map be recorded by finitely many numbers? Once bases are fixed, record the coordinates of each basis image in one column.\end{qbox}
\begin{prob}\begin{definition*}[Notes Definition 16; LADR 3.29]An $m\times n$ matrix $A=(a_{ij})$ has $m$ rows and $n$ columns, with $a_{ij}\in F$ in row $i$, column $j$. The vector space $F^{m,n}$ uses entrywise addition and scaling. Its vector-space axioms hold entry by entry, by the field axioms. We write coordinate vectors as columns. The identity matrix $I_n$ has ones on its diagonal and zeros elsewhere; its columns are $e_1,\ldots,e_n$.\end{definition*}\end{prob}
\begin{proposition*}[Notes Chapter 6; LADR 3.40]The matrices $E_{ij}$ with a single $1$ at position $(i,j)$ form a basis of $F^{m,n}$; hence its dimension is $mn$.\end{proposition*}
\begin{motive}\ii[Want.] List the independent entries of an array.\ii[Know.] Equality of matrices means entrywise equality.\ii[Key step.] Use one matrix for each entry.\end{motive}
\begin{prob}\begin{proof}\begin{steps}\ii Every $A$ equals $\sum_{i=1}^m\sum_{j=1}^n a_{ij}E_{ij}$, proving spanning.\ii The $(i,j)$ entry of this sum is exactly $a_{ij}$; a zero sum forces every coefficient zero. There are $mn$ basis matrices.\end{steps}\end{proof}\end{prob}
\begin{example*}[Notes Examples 26--27]\begin{steps}\ii Entrywise operations give
\[\begin{pmatrix}1+i&3+2i&4\\5&i&2+i\end{pmatrix}+\begin{pmatrix}4-2i&5&1\\3&2&i\end{pmatrix}=\begin{pmatrix}5-i&8+2i&5\\8&2+i&2+2i\end{pmatrix}.\]
\ii Scaling gives $\displaystyle(1+i)\begin{pmatrix}1+i&3+2i&4\\5&i&2+i\end{pmatrix}=\begin{pmatrix}2i&1+5i&4+4i\\5+5i&-1+i&1+3i\end{pmatrix}$.
\ii For $F^{2,2}$ the ordered basis is $\left(\begin{smallmatrix}1&0\\0&0\end{smallmatrix}\right),\left(\begin{smallmatrix}0&1\\0&0\end{smallmatrix}\right),\left(\begin{smallmatrix}0&0\\1&0\end{smallmatrix}\right),\left(\begin{smallmatrix}0&0\\0&1\end{smallmatrix}\right)$.\end{steps}\end{example*}
\begin{prob}\begin{definition*}[Notes Definitions 18, 17; LADR 3.73, 3.31]For an ordered basis $B=(v_1,\ldots,v_n)$, the coordinate column $[v]_B=(a_1,\ldots,a_n)^{\mathsf T}$ is defined by $v=\sum a_jv_j$. Here the superscript $\mathsf T$ converts a displayed row into a column. For a codomain basis $C=(w_1,\ldots,w_m)$, define $M_C^B(T)$ to have column $j$ equal to $[Tv_j]_C$. Equivalently $Tv_j=\sum_i a_{ij}w_i$. The ordering of both bases matters.\end{definition*}\end{prob}
\begin{example*}[Notes Example 29]\begin{steps}\ii For $p=1+2x+3x^2$, the coordinates in $(1,x,x^2)$ are $(1,2,3)^{\mathsf T}$.
\ii The list $B=(1+x,1+x^2,x+x^2)$ is independent: $a+b=0$, $a+c=0$, $b+c=0$ force $2a=0$, hence $a=b=c=0$ over $\mathbb R$. Its length $3$ gives a basis of $\mathcal P_2(\mathbb R)$.
\ii Since $p=0(1+x)+(1+x^2)+2(x+x^2)$, $[p]_B=(0,1,2)^{\mathsf T}$.\end{steps}\end{example*}
\method{Build a representation matrix}\begin{steps}\ii Fix the ordered domain and codomain bases. Apply $T$ to the first domain vector and express the result in the codomain basis; put those coefficients in column 1.\ii Repeat for each domain vector. For $P,R_\theta,N$ in the standard basis this gives\[M(P)=\begin{pmatrix}0&0\\0&1\end{pmatrix},\quad M(R_\theta)=\begin{pmatrix}\cos\theta&-\sin\theta\\\sin\theta&\cos\theta\end{pmatrix},\quad M(N)=\begin{pmatrix}0&1\\0&0\end{pmatrix}.\]\ii The column count is the domain dimension; the row count is the codomain dimension.\end{steps}
\begin{example*}[Notes Example 28]Let $T:\mathcal P_3(\mathbb C)\to\mathbb C^4$ send $p$ to $(p(1),p(i),p(-1),p(-i))$.
\begin{steps}\ii Evaluation is linear: $(ap+bq)(z)=ap(z)+bq(z)$ at each of the four points.
\ii Apply $T$ to $1,x,x^2,x^3$, successively; powers of $i$ give columns of the standard-basis matrix
\[A=M(T)=\begin{pmatrix}1&1&1&1\\1&i&-1&-i\\1&-1&1&-1\\1&-i&-1&i\end{pmatrix}.\]\end{steps}\end{example*}
\begin{proposition*}[Notes Proposition 26; LADR 3.35, 3.38, 3.71--3.72]For fixed bases of finite-dimensional $V,W$, $T\mapsto M_C^B(T)$ is a linear bijection $L(V,W)\to F^{m,n}$. Consequently $\dim L(V,W)=(\dim V)(\dim W)$.\end{proposition*}
\begin{motive}\ii[Want.] Replace maps by arrays without losing information.\ii[Know.] Proposition 21 allows arbitrary basis images, uniquely.\ii[Key step.] Read a proposed array column by column as prescribed images.\end{motive}
\begin{prob}\begin{proof}\begin{steps}\ii Every matrix gives unique vectors $Tv_j$ by using its columns as $C$-coordinates. Proposition 21 gives exactly one map with those values, proving bijection.\ii The coordinates of $(S+T)v_j$ and $aTv_j$ are the column sums and scalar multiples. Thus $M(S+T)=M(S)+M(T)$ and $M(aT)=aM(T)$.\ii The inverse images of the $mn$ matrices $E_{ij}$ form a basis: expand a matrix in them for spanning, and apply the injective correspondence to a zero relation for independence.\end{steps}\end{proof}\end{prob}
\begin{shortcuts}\ii Apply a map only to basis vectors; Proposition 21 supplies all other values.\ii Use polynomial coefficient comparison to find coordinates; a coefficient list is a coordinate vector only for the specified basis.\end{shortcuts}\begin{trapbox}\ii Columns describe images, not rows: $N(e_2)=e_1$ places its $1$ in row 1, column 2.\ii The same polynomial has coordinates $(1,2,3)$ and $(0,1,2)$ above; changing a basis changes its coordinates.\end{trapbox}
\section{Matrix products and ranks}\begin{qbox}Why does row-times-column multiplication represent composition? Each intermediate coordinate contributes its coefficient times the corresponding next image.\end{qbox}
\begin{prob}\begin{definition*}[Notes Definitions 19--20; LADR 3.41, 3.76]For $A\in F^{m,n}$ and $x\in F^n$, set $(Ax)_i=\sum_{j=1}^n a_{ij}x_j$. For $B\in F^{n,r}$, define the $m\times r$ product by $(AB)_{ik}=\sum_{j=1}^n a_{ij}b_{jk}$. Both definitions depend only on the entries. The notation $A_{:j}$ means column $j$ and $A_{i:}$ means row $i$. Distributivity of finite sums shows that $x\mapsto Ax$ is linear (the matrix-map part of Notes Example 21).\end{definition*}\end{prob}
\begin{proposition*}[Notes coordinate/product formulas; LADR 3.43, 3.46, 3.48, 3.50--3.51, 3.76]For compatible bases, $[Tv]_C=M_C^B(T)[v]_B$ and $M_D^B(ST)=M_D^C(S)M_C^B(T)$. The $k$th product column is $A$ times column $k$ of $B$. Also $Ax=\sum_jx_jA_{:j}$. Each product column is a combination of columns of $A$; each product row is a combination of rows of $B$.\end{proposition*}
\begin{motive}\ii[Want.] Translate linear-map operations into array operations.\ii[Know.] Basis images are matrix columns, and linearity distributes coefficients.\ii[Key step.] Expand one input and then collect each output basis coefficient.\end{motive}
\begin{prob}\begin{proof}\begin{steps}\ii If $v=\sum_jx_jv_j$ and $Tv_j=\sum_i a_{ij}w_i$, then $Tv=\sum_i(\sum_j a_{ij}x_j)w_i$. This proves the coordinate identity and the column combination formula.\ii Apply this identity to $STv_k=S(Tv_k)$. Its coordinate column is $M(S)$ times column $k$ of $M(T)$, yielding the product formula.\ii The entry formula gives $(AB)_{:k}=\sum_jb_{jk}A_{:j}$ and $(AB)_{i:}=\sum_ja_{ij}B_{j:}$. Here $:$ means all entries of that column or row.\end{steps}\end{proof}\end{prob}
\begin{example*}[Notes Examples 30--31]\begin{steps}\ii Row-times-column multiplication gives\[\begin{pmatrix}1&2&3\\2&3&4\end{pmatrix}\begin{pmatrix}5\\7\\9\end{pmatrix}=\begin{pmatrix}5+14+27\\10+21+36\end{pmatrix}=\begin{pmatrix}46\\67\end{pmatrix}.\]
\ii For the evaluation matrix $A$ in Example 28, $A(1,2,1,0)^{\mathsf T}=(4,2i,0,-2i)^{\mathsf T}$. Directly evaluating $1+2x+x^2$ at $1,i,-1,-i$ gives the same four values.\end{steps}\end{example*}
\begin{example*}[Notes Example 32]Let $T:\mathcal P_3(\mathbb R)\to\mathcal P_2(\mathbb R)$ be differentiation and $S:\mathcal P_2(\mathbb R)\to\mathcal P_3(\mathbb R)$ send $p$ to $\int_0^x p(t)\,dt$.
\begin{steps}\ii The power rule gives $M(T)=\begin{pmatrix}0&1&0&0\\0&0&2&0\\0&0&0&3\end{pmatrix}$ and $M(S)=\begin{pmatrix}0&0&0\\1&0&0\\0&1/2&0\\0&0&1/3\end{pmatrix}$.
\ii Their product $M(ST)=M(S)M(T)=\operatorname{diag}(0,1,1,1)$, where $\operatorname{diag}$ means the displayed diagonal entries and zeros elsewhere: $STp=p-p(0)$.
\ii Reversing the order gives $M(TS)=I_3$ and $TSp=p$. The two products even have different sizes.\end{steps}\end{example*}
\begin{proposition*}[Notes Proposition 27; LADR 3.8, 3.43]Matrix products are associative and distribute on both sides: $(AB)C=A(BC)$, $A(B+C)=AB+AC$, $(A+B)C=AC+BC$, whenever the sizes permit the operations. The identity matrix $I_n$ has ones on its diagonal and zeros elsewhere; $AI_n=A$ and $I_mA=A$ for $A\in F^{m,n}$.\end{proposition*}
\begin{motive}\ii[Want.] Use algebraic laws without repeatedly expanding all entries.\ii[Know.] Matrices represent maps between coordinate spaces, and products represent composition.\ii[Key step.] Transfer the proved laws of composition through the matrix correspondence.\end{motive}
\begin{prob}\begin{proof}\begin{steps}\ii Associate each matrix with its map $x\mapsto Ax$. The composition laws of Proposition 19 give associativity and both distributive laws for these maps.\ii Equality of their matrices gives the stated array identities. The identity map sends each $e_j$ to itself, so its matrix is $I_n$; the identity laws follow.\end{steps}\end{proof}\end{prob}
\begin{prob}\begin{definition*}[LADR 3.52, 3.54, 3.58]The column rank of $A$ is the dimension of the span of its columns; its row rank is the dimension of the span of its rows (viewed as vectors). The transpose $A^{\mathsf T}$ exchanges rows and columns: $(A^{\mathsf T})_{ji}=a_{ij}$. We prove the two ranks agree and call their common value the rank of $A$.\end{definition*}\end{prob}
\begin{proposition*}[LADR 3.54: transpose algebra]
For compatible matrices, $(A+B)^{\mathsf T}=A^{\mathsf T}+B^{\mathsf T}$, $(aA)^{\mathsf T}=aA^{\mathsf T}$ and $(AB)^{\mathsf T}=B^{\mathsf T}A^{\mathsf T}$.
\end{proposition*}
\begin{motive}\ii[Want.] Transpose expressions correctly.\ii[Know.] Transposition exchanges the two indices.\ii[Key step.] Compare the $(i,j)$ entries; the product order reverses because the intermediate index changes position.\end{motive}
\begin{prob}\begin{proof}\begin{steps}
\ii The $(i,j)$ entries in the sum identity are both $a_{ji}+b_{ji}$; those in the scalar identity are both $aa_{ji}$ (the first $a$ is the scalar).
\ii The $(i,j)$ entry of $(AB)^{\mathsf T}$ is $\sum_k a_{jk}b_{ki}=\sum_k(B^{\mathsf T})_{ik}(A^{\mathsf T})_{kj}$, the corresponding entry of $B^{\mathsf T}A^{\mathsf T}$.
\end{steps}\end{proof}\end{prob}
\begin{proposition*}[LADR 3.56--3.57, 3.78]If $A\in F^{m,n}$ has positive column rank $c$, then $A=CR$ where the $c$ columns of $C\in F^{m,c}$ form a basis of the column span, and $R\in F^{c,n}$ contains the coefficients of the columns of $A$ in that basis. Row rank equals column rank. Also $\operatorname{rank}M(T)=\dim\im T$.\end{proposition*}
\begin{motive}\ii[Want.] Count independent rows and columns without separate full computations.\ii[Know.] Every column expands in a chosen column basis; rows of a product are combinations of rows of its second factor.\ii[Key step.] Factor through the column-basis coefficients, then apply the same argument to the transpose.\end{motive}
\begin{prob}\begin{proof}\begin{steps}\ii Choose independent columns spanning all columns of $A$ and put them in $C$. Put their expansion coefficients in $R$. The column formula for products gives $A=CR$.\ii Every row of $A=CR$ is a combination of the $c$ rows of $R$. Thus row rank $\le c=$ column rank. For the zero matrix both ranks are zero.\ii Apply the inequality to $A^{\mathsf T}$. Its rows are the columns of $A$ and conversely, so column rank of $A\le$ row rank of $A$. Equality follows.\ii The coordinate map $w\mapsto[w]_C$ is linear and bijective by unique basis coefficients. The range is spanned by $Tv_j$; their coordinate columns are the columns of $M(T)$. This bijection preserves spanning and independence, so their dimensions agree.\end{steps}\end{proof}\end{prob}
\method{Find a column--row factorization}\begin{steps}\ii Let $A=\begin{pmatrix}1&0&2\\0&1&3\\1&1&5\end{pmatrix}$. Its first two columns are independent by their first two entries; column 3 is twice column 1 plus three times column 2.\ii Therefore $C=\begin{pmatrix}1&0\\0&1\\1&1\end{pmatrix}$ and $R=\begin{pmatrix}1&0&2\\0&1&3\end{pmatrix}$ give $A=CR$, rank $2$. The row relation is row 3 = row 1 + row 2.\ii Solving $Ax=0$ gives $x_1=-2x_3,x_2=-3x_3$, so the kernel of the matrix map is $\operatorname{span}((-2,-3,1))$, agreeing with $3=1+2$.\end{steps}
\begin{shortcuts}\ii For $AB$, check inner dimensions first; the result has the outer dimensions.\ii Use a product column $A(B_{:j})$ when only that column is needed.\ii A column basis computes the rank and simultaneously gives a basis of the coordinate range.\end{shortcuts}\begin{trapbox}\ii $AB=0$ does not imply either factor is zero: $N^2=0$ but $N\ne0$.\ii Matrix multiplication is not entrywise: $N^2=0$, whereas squaring each entry of $N$ returns $N$.\ii Row rank equals column rank, but the row and column spaces usually live in different spaces: a $2\times3$ matrix has rows in $F^3$ and columns in $F^2$.\end{trapbox}
\section{Permutation and stochastic matrices}\begin{qbox}How does a matrix update a distribution? Column $j$ tells us the next distribution when the present state is certainly $j$.\end{qbox}
\begin{prob}\begin{definition*}[Notes Chapter 6]A permutation $\sigma$ of $\{1,\ldots,n\}$ is a bijection of this set. Its permutation matrix is defined by $S_\sigma e_j=e_{\sigma(j)}$: every row and column contains exactly one $1$, all other entries being zero. A probability vector $p\in\mathbb R^n$ has nonnegative entries summing to $1$. A column-stochastic matrix has nonnegative entries and each column sums to $1$.\end{definition*}\end{prob}
\begin{proposition*}[Notes Chapter 6: repeated transitions]$S_\sigma^k=S_{\sigma^k}$ for all integers $k\ge0$. A column-stochastic matrix $P$ sends probability vectors to probability vectors. A time-homogeneous finite Markov chain uses the same transition probabilities at every step and has next-state probabilities depending only on its present state; its distributions satisfy $p_k=P^kp_0$.\end{proposition*}
\begin{motive}\ii[Want.] Describe repeated transitions and preserve total probability.\ii[Know.] Composition is determined on a basis, and matrix products sum over the intermediate state.\ii[Key step.] Track each basis state or sum the coordinates of a distribution.\end{motive}
\begin{prob}\begin{proof}\begin{steps}\ii On $e_j$, $S_\sigma^ke_j=e_{\sigma^k(j)}$ by induction, proving the matrix identity by basis determination.\ii If $p$ is a probability vector, $(Pp)_i=\sum_jP_{ij}p_j\ge0$, and $\sum_i(Pp)_i=\sum_j(\sum_iP_{ij})p_j=\sum_jp_j=1$.\ii Here $P_{ij}$ is the probability of moving to $i$ given present state $j$. Summing the mutually exclusive present states gives $p_{k+1}=Pp_k$; induction gives $p_k=P^kp_0$.\end{steps}\end{proof}\end{prob}
\begin{example*}[Notes Example 33]Take $P=\begin{pmatrix}1/2&1/4&0\\1/2&1/2&1/2\\0&1/4&1/2\end{pmatrix}$ and $p_0=(1,0,0)^{\mathsf T}$.
\begin{steps}\ii Every column sums to $1$ and is nonnegative. From state 1, the next state is 1 or 2 with probability $1/2$ each; column 2 gives probabilities $1/4,1/2,1/4$ from state 2.
\ii Repeated multiplication gives $p_1=(1/2,1/2,0)^{\mathsf T}$, $p_2=(3/8,1/2,1/8)^{\mathsf T}$, $p_3=(5/16,1/2,3/16)^{\mathsf T}$.
\ii Direct multiplication gives $P\pi=\pi$ for $\pi=(1/4,1/2,1/4)^{\mathsf T}$. Thus once this distribution is reached, it remains unchanged.\end{steps}\end{example*}
\method{Switch between probability conventions}\begin{steps}\ii Transposing the coordinate equation $(Pp)_i=\sum_jP_{ij}p_j$ gives $(Pp)^{\mathsf T}=p^{\mathsf T}P^{\mathsf T}$.\ii Thus the row-vector convention uses the row-stochastic matrix $P^{\mathsf T}$. Permutation matrices are column-stochastic and describe deterministic transitions $j\mapsto\sigma(j)$.\end{steps}
\begin{shortcuts}\ii Check stochasticity column by column; preservation of total probability then follows without testing individual distributions.\ii To verify an unchanged distribution, check $P\pi=\pi$ and that $\pi$ is a probability vector.\end{shortcuts}\begin{trapbox}\ii Row sums equal to $1$ are insufficient for our column convention: $\left(\begin{smallmatrix}1&0\\1&0\end{smallmatrix}\right)e_1=(1,1)^{\mathsf T}$ has total $2$.\ii An unchanged distribution need not be approached from other starting distributions: $I_np=p$ for every $p$.\end{trapbox}
\section{Problem set}\begin{qbox}Use basis images, kernel/range, or dimensions before expanding long equations. Each problem below comes from HW3 or the course practice midterm.\end{qbox}
\begin{prob}\textbf{HW3 P1.} Let $T:V\to W$ be linear with $V$ finite-dimensional, and $U\subseteq W$ a subspace. (a) Prove $T^{-1}(U)=\{v:Tv\in U\}$ is a subspace. (b) Prove $\dim T^{-1}(U)=\dim\ker T+\dim(U\cap\im T)$. (c) For differentiation $D:\mathcal P_4(\mathbb R)\to\mathcal P_4(\mathbb R)$, find the dimension of the polynomials satisfying $p'(x)=ax^2+b$ for some $a,b\in\mathbb R$.\\\emph{Hint:} Apply rank--nullity to $T:T^{-1}(U)\to U\cap\im T$.\end{prob}
\begin{prob}\textbf{HW3 P2.} (a) Find $T:\mathbb R^4\to\mathbb R^4$ linear with $\ker T=\im T$. (b) Prove no such map exists on $\mathbb R^5$.\\\emph{Hint:} Use two coordinates as outputs; apply Proposition 24 to equal kernel and range dimensions.\end{prob}
\begin{prob}\textbf{HW3 P3.} For $U=\{p\in\mathcal P_4(\mathbb R):p(1)=p'(1)=0\}$, define $S:\mathcal P_4(\mathbb R)\to\mathbb R^2$ with kernel $U$. Use rank--nullity for $\dim U$ and compare with HW2.\\\emph{Hint:} Use the two constraints as outputs; test $1$ and $x-1$ for surjectivity.\end{prob}
\begin{prob}\textbf{HW3 P4.} For linear $T:V\to W$, (a) prove a dependent list has a dependent image list, and (b) show an independent list need not have an independent image list.\\\emph{Hint:} Apply $T$ to a nontrivial relation; use the zero map for (b).\end{prob}
\begin{prob}\textbf{HW3 P5.} Let $T:\mathcal P_3(\mathbb R)\to\mathcal P_4(\mathbb R)$ be $(Tp)(x)=x^2p'(x)+p(x)$. (a) Find its matrix in the standard monomial bases. (b) Find kernel and range and their dimensions. (c) Explain directly why the range is the span of the representing columns, under monomial coordinates.\\\emph{Hint:} Apply $T$ to $1,x,x^2,x^3$ and use Proposition 21.\end{prob}
\begin{prob}\textbf{HW3 P6.} Over $\mathbb F_2$, solve $x+y+z+w=1$, $x+z=0$, listing every solution.\\\emph{Hint:} Solve for $z,w$ in terms of $x,y$ using addition modulo 2.\end{prob}
\begin{prob}\textbf{HW3 P7 (optional).} Let $A$ be an $n\times m$ matrix over $\mathbb F_2$ and suppose $Ax=b$ is consistent (has at least one solution). Express the number of solutions using $m$ and $\operatorname{rank}A$.\\\emph{Hint:} Fix one solution $x_0$; every other solution differs by an element of $\ker A$. Count basis coefficients.\end{prob}
\begin{prob}\textbf{Exam practice (practice-midterm1.pdf, Question 2).} Let $V,W$ be finite-dimensional and $T:V\to W$ injective and linear. (a) For a basis $v_1,\ldots,v_n$ of $V$, prove $Tv_1,\ldots,Tv_n$ is independent. (b) Prove there exists a linear map $S:W\to V$ with $ST=I_V$. Such an $S$ is called a left inverse of $T$.\\\emph{Hint:} Use injectivity for (a), then extend the image list to a basis and prescribe $S$ there by Proposition 21.\end{prob}\end{document}
